\pagenumbering{roman}
\setcounter{page}{1}

\selecthungarian

%----------------------------------------------------------------------------
% Abstract in Hungarian
%----------------------------------------------------------------------------
\chapter*{Kivonat}\addcontentsline{toc}{chapter}{Kivonat}

A pénzügyi szektor digitális transzformációjának egyik legnagyobb gátja a mai napig a manuális, papír- vagy PDF-alapú dokumentumfeldolgozás. A vállalati gyakorlatban a számlák jelentős része strukturálatlan formában érkezik, melyek feldolgozása a különböző rendszerek közötti manuális adatátvitellel („swivel-chair integration”) történik. Ez a folyamat nemcsak erőforrás-igényes és lassú, de magas hibakockázatot és átláthatatlanságot is eredményez a beszerzéstől a fizetésig (Purchase-to-Pay) tartó láncban.

Szakdolgozatomban ezen problémára kínálok megoldást a Hiperautomatizálás (Hyperautomation) eszköztárának alkalmazásával. A dolgozat célja egy olyan integrált, „end-to-end” számlafeldolgozó prototípus tervezése és megvalósítása az SAP Business Technology Platform (BTP) felhőalapú környezetében, amely minimalizálja az emberi beavatkozást, ugyanakkor biztosítja a folyamat kontrollálhatóságát.

A megvalósítás során az SAP Build Process Automation (SBPA) technológiát alkalmaztam, ötvözve a robotikus folyamatautomatizálást (RPA) és a mesterséges intelligenciát. A rendszer bemenetét egy saját fejlesztésű, determinisztikus Python szkripttel előállított szintetikus számlaállomány biztosítja, amely lehetővé tette a különböző hibatípusok (hiányos adat, formátumhiba) szisztematikus tesztelését. A strukturálatlan adatok kinyerését az SAP Document Information Extraction (DOX) szolgáltatás legújabb, Generative AI alapú modelljével valósítottam meg. A rendszer architektúrájában kiemelt szerepet kapott az egyedi JavaScript kódokkal vezérelt API kommunikáció és adattranszformáció, amely a technikai robot (Bot) és az üzleti folyamat (Process) közötti hidat képezi. A felhasználói interakciókat és a jóváhagyási döntéseket („Human-in-the-loop”) az SAP Build Work Zone egységesített felületén integráltam.

A prototípus tesztelése három különböző forgatókönyv alapján igazolta, hogy a megoldás képes a számlák teljesen automatikus („touchless”) feldolgozására, valamint a kivételes esetek intelligens kezelésére. Az eredmények azt mutatják, hogy az automatizált folyamat a manuális gyakorlathoz képest radikálisan csökkenti az átfutási időt és az adatrögzítési hibák számát, miközben valós idejű átláthatóságot biztosít a pénzügyi folyamatokban.

\vfill
\selectenglish


%----------------------------------------------------------------------------
% Abstract in English
%----------------------------------------------------------------------------
\chapter*{Abstract}\addcontentsline{toc}{chapter}{Abstract}

One of the biggest obstacles to digital transformation in the financial sector remains manual, paper- or PDF-based document processing. In corporate practice, a significant portion of invoices arrives in unstructured formats, requiring manual data transfer between systems ("swivel-chair integration"). This process is not only resource-intensive and slow but also results in high error risks and a lack of transparency in the Purchase-to-Pay chain.

In my thesis, I propose a solution to this problem by applying the toolkit of Hyperautomation. The aim of the thesis is to design and implement an integrated, end-to-end invoice processing prototype within the cloud-based environment of the SAP Business Technology Platform (BTP), minimizing human intervention while ensuring process control.

For the implementation, I utilized SAP Build Process Automation (SBPA) technology, combining Robotic Process Automation (RPA) with Artificial Intelligence. The system input is provided by a synthetic invoice dataset generated by a self-developed deterministic Python script, which allowed for systematic testing of various error types (missing data, format errors). The extraction of unstructured data was achieved using the latest Generative AI-based model of the SAP Document Information Extraction (DOX) service. A key component of the system architecture is API communication and data transformation driven by custom JavaScript codes, bridging the technical bot and the business process. User interactions and approval decisions ("Human-in-the-loop") were integrated into the unified interface of SAP Build Work Zone.

Testing the prototype across three different scenarios confirmed that the solution is capable of fully automated ("touchless") invoice processing as well as intelligent handling of exception cases. The results demonstrate that the automated process radically reduces turnaround time and data entry errors compared to manual practices, while providing real-time transparency in financial processes.

\vfill
\cleardoublepage

\selectthesislanguage

\newcounter{romanPage}
\setcounter{romanPage}{\value{page}}
\stepcounter{romanPage}